% Contexto e instrucciones literales de las páginas 1 y 2 de EXA-0053.
% Fuente: Selectividad/sel_2013_matematicas_aplicadas/septiembre/titular/EXAMEN 2-ANDALUCÍA.pdf
% SHA-256: f8633bbd91d33bee02c2cc582c33b7ffeb3643f312b810ad0bb8c5d154d0721d
% No se reproduce el logotipo ni el marco de la cabecera.
\section*{Contexto del examen original}
\textbf{UNIVERSIDADES DE ANDALUCÍA}\par
\textbf{PRUEBA DE ACCESO A LA UNIVERSIDAD}\par
CURSO 2012-2013\par
\textbf{MATEMÁTICAS APLICADAS A LAS CIENCIAS SOCIALES II}

\textbf{Instrucciones:}
\begin{enumerate}[label=\textbf{\alph*)}]
  \item Duración: 1 hora y 30 minutos.
  \item Elija una de las dos opciones propuestas y conteste los ejercicios de la opción elegida.
  \item En cada ejercicio, parte o apartado se indica la puntuación máxima que le corresponde.
  \item Se permitirá el uso de calculadoras que no sean programables, gráficas ni con capacidad para almacenar o transmitir datos.
  \item Si obtiene resultados directamente con la calculadora, explique con detalle los pasos necesarios para su obtención sin su ayuda. Justifique las respuestas.
\end{enumerate}
\noindent\textbf{Nota de la colección:} Se reúnen ambas opciones para consulta. La regla anterior pertenece al examen original. Curso 2012-2013 impreso. Septiembre y titular sólo de carpeta,no impresos. Sin letra de reserva;modelo no confirmado;2 no acredita modelo. Examen y criterios sin alias exactos. Tabla normal SRC-B65557E07CB9 asociada por carpeta,sin alias exactos en el inventario.
\par\medskip
